%%%%%%%%%%%%%%%%%%%%%%%%%%%%%%%%%%%%%%%%%%%%%%%%%%%%%%%%%
% Presentasi Beamer untuk RPS Optimisasi Teknik Elektro
% Minggu ke-1: Pengantar Optimisasi & Pemodelan Matematis
% Program Studi S1 Teknik Elektro - Universitas Bengkulu (2026)
%%%%%%%%%%%%%%%%%%%%%%%%%%%%%%%%%%%%%%%%%%%%%%%%%%%%%%%%%

\documentclass[10pt,aspectratio=169]{beamer}

% --- TEMA DAN WARNA ---
\usetheme{Madrid}
\usecolortheme{seahorse}
\setbeamertemplate{navigation symbols}{}

% Custom UNIB Colors
\definecolor{unibblue}{RGB}{0, 51, 102}
\definecolor{unibgold}{RGB}{212, 175, 55}
\definecolor{darkgreen}{RGB}{34, 139, 34}

\setbeamercolor{structure}{fg=unibblue}
\setbeamercolor{palette primary}{bg=unibblue,fg=white}
\setbeamercolor{palette secondary}{bg=unibblue!80,fg=white}
\setbeamercolor{palette tertiary}{bg=unibblue!60,fg=white}
\setbeamercolor{titlelike}{parent=palette primary,fg=white}

% --- PACKAGES ---
\usepackage[utf8]{inputenc}
\usepackage[indonesian]{babel}
\usepackage{amsmath, amssymb, amsfonts}
\usepackage{graphicx}
\usepackage{booktabs}
\usepackage{tikz}
\usepackage{pgfplots}
\pgfplotsset{compat=1.17}
\usepackage{caption}
\usepackage{listings}
\usepackage{array}
\usepackage{tcolorbox}
\tcbuselibrary{skins,breakable}

% Listings setup for Python/Julia
\lstset{
  language=Python,
  basicstyle=\ttfamily\tiny,
  keywordstyle=\color{unibblue}\bfseries,
  commentstyle=\color{darkgreen}\itshape,
  stringstyle=\color{red!70!black},
  numbers=left,
  numberstyle=\tiny\color{gray},
  numbersep=4pt,
  breaklines=true,
  frame=single,
  rulecolor=\color{unibblue!30},
  tabsize=4
}

% --- TITLE PAGE INFORMATION ---
\title[Optimisasi TE: Pertemuan 1]{Optimisasi untuk Teknik Elektro (TKE-41XX)}
\subtitle{Pertemuan 1: Pengantar Optimisasi dan Pemodelan Matematis}
\author[N. Daratha]{Ir. Novalio Daratha S.T., M.Sc., Ph.D.}
\institute[Teknik Elektro UNIB]{Program Studi S1 Teknik Elektro\\Fakultas Teknik --- Universitas Bengkulu}
\date{Semester Ganjil 2026}
% --- CUSTOM COMMANDS ---
\newcommand{\R}{\mathbb{R}}

\begin{document}

% --- SLIDE 1: JUDUL ---
\begin{frame}
    \titlepage
\end{frame}

% --- SLIDE 2: AGENDA PERTEMUAN ---
\begin{frame}{Agenda Pembelajaran Minggu ke-1}
    \begin{tcolorbox}[colback=unibblue!5,colframe=unibblue,title=\textbf{Tujuan Pembelajaran (CPMK-1 dan CPMK-2)}]
        Mahasiswa mampu menjelaskan konsep dasar optimisasi, mengklasifikasikan masalah rekayasa, dan menerjemahkan masalah praktis teknik elektro ke dalam bentuk pemodelan matematis standar.
    \end{tcolorbox}

    \vspace{0.3cm}
    \begin{enumerate}
        \item \textbf{Pengantar Optimisasi}: Peranan optimisasi dalam sistem tenaga, telekomunikasi, dan kontrol.
        \item \textbf{Anatomi Model Optimisasi}: Variabel keputusan, fungsi tujuan, dan kendala (*constraints*).
        \item \textbf{Klasifikasi Masalah Optimisasi}: LP, NLP, Cembung vs Non-Cembung, MIP, dan Multi-Objektif.
        \item \textbf{Studi Kasus Sederhana}: Pengalokasian daya 2 pembangkit (*Economic Dispatch* sederhana).
        \item \textbf{Pengenalan Perangkat Komputasi}: Python (SciPy) dan Julia (JuMP).
    \end{enumerate}
\end{frame}

% --- SECTION 1 ---
\section{Pengantar Optimisasi}

\begin{frame}{Apa itu Optimisasi?}
    \begin{columns}
        \column{0.55\textwidth}
        \begin{block}{Definisi Utama}
            \textbf{Optimisasi} adalah seni dan sains untuk mencapai **hasil terbaik** (biaya terendah, efisiensi tertinggi, atau keuntungan maksimal) dari suatu sistem rekayasa di bawah keterbatasan (*kendala*) sumber daya yang ada.
        \end{block}

        \vspace{0.2cm}
        \begin{alertblock}{Mengapa Penting untuk Insinyur Elektro?}
          \begin{itemize}
              \item Desain jaringan listrik hemat biaya dan rugi-rugi minimal.
              \item Alokasi frekuensi dan penempatan antena seluler.
              \item Penalaan parameter kontroler PID presisi tinggi.
          \end{itemize}
        \end{alertblock}

        \column{0.41\textwidth}
        \begin{tcolorbox}[colback=unibgold!10,colframe=unibgold!80!black,title=\textbf{3 Elemen Kunci}]
            \begin{enumerate}
                \item \textbf{Apa yang dicari?} (Variabel $\mathbf{x}$)
                \item \textbf{Apa ukurannya?} (Tujuan $f_0(\mathbf{x})$)
                \item \textbf{Apa batasannya?} (Kendala $g_i, h_j$)
            \end{enumerate}
        \end{tcolorbox}
    \end{columns}
\end{frame}

% --- SECTION 2 ---
\section{Anatomi Model Optimisasi}

\begin{frame}{Formulasi Standar Masalah Optimisasi}
    Secara matematis, bentuk umum optimisasi terkendala dituliskan sebagai:
    \begin{block}{Formulasi Standar Minimisasi}
        \[
            \min_{\mathbf{x} \in \R^n} f_0(\mathbf{x})
        \]
        \[
            \text{dengan kendala:} \quad
            \begin{cases}
                g_i(\mathbf{x}) \le 0, & i = 1, 2, \dots, m \quad (\text{Pertidaksamaan}) \\
                h_j(\mathbf{x}) = 0, & j = 1, 2, \dots, p \quad (\text{Persamaan}) \\
                x_{k,\min} \le x_k \le x_{k,\max}, & k = 1, 2, \dots, n \quad (\text{Batas Domain})
            \end{cases}
        \]
    \end{block}

    \begin{tcolorbox}[colback=white,colframe=unibblue,title=\textbf{Catatan Penting}]
        Masalah maksimisasi $\max f_0(\mathbf{x})$ setara dengan minimisasi $-\min [-f_0(\mathbf{x})]$.
    \end{tcolorbox}
\end{frame}

\begin{frame}{Klasifikasi Masalah Optimisasi}
    \begin{table}
        \centering
        \small
        \begin{tabular}{p{3.5cm}p{5.5cm}p{4cm}}
            \toprule
            \textbf{Kategori} & \textbf{Karakteristik Matematis} & \textbf{Contoh Teknik Elektro} \\
            \midrule
            Linear Programming (LP) & Fungsi tujuan dan kendala linier & Alokasi beban trafo \\
            Non-Linear (NLP) & Mengandung fungsi non-linier & Economic Dispatch \\
            Mixed-Integer (MIP) & Variabel bernilai biner (0, 1) & Penempatan antena \\
            Convex Optimization & Himpunan dan fungsi cembung & Optimum global tunggal \\
            \bottomrule
        \end{tabular}
    \end{table}
\end{frame}

% --- SECTION 3 ---
\section{Studi Kasus Sederhana}

\begin{frame}{Studi Kasus 1: Pembagian Beban 2 Generator}
    \begin{block}{Deskripsi Masalah}
        Dua generator termal melayani total beban $P_D = 400\text{ MW}$.
        \begin{itemize}
            \item Generator 1: $F_1(P_1) = 0.002 P_1^2 + 8.0 P_1 + 100 \quad [\text{\$/jam}], \quad 50 \le P_1 \le 300\text{ MW}$
            \item Generator 2: $F_2(P_2) = 0.003 P_2^2 + 7.5 P_2 + 150 \quad [\text{\$/jam}], \quad 50 \le P_2 \le 250\text{ MW}$
        \end{itemize}
    \end{block}

    \begin{tcolorbox}[colback=white,colframe=darkgreen,title=\textbf{Model Optimisasi Matematis}]
        \[
            \min_{P_1, P_2} F_T = (0.002 P_1^2 + 8.0 P_1 + 100) + (0.003 P_2^2 + 7.5 P_2 + 150)
        \]
        \[
            \text{s.t.} \quad P_1 + P_2 = 400, \quad 50 \le P_1 \le 300, \quad 50 \le P_2 \le 250
        \]
    \end{tcolorbox}
\end{frame}

\begin{frame}[fragile]{Implementasi Pertama dengan Python (SciPy)}
\begin{lstlisting}
import numpy as np
from scipy.optimize import minimize

# Fungsi Tujuan (Total Biaya)
def objective(P):
    P1, P2 = P[0], P[1]
    cost1 = 0.002 * P1**2 + 8.0 * P1 + 100
    cost2 = 0.003 * P2**2 + 7.5 * P2 + 150
    return cost1 + cost2

# Kendala Persamaan (P1 + P2 = 400)
constraints = ({'type': 'eq', 'fun': lambda P: P[0] + P[1] - 400})

# Batas Operasi (Bounds)
bounds = [(50, 300), (50, 250)]

# Eksekusi Minimisasi
res = minimize(objective, [200, 200], method='SLSQP', bounds=bounds, constraints=constraints)
print(f"P1 Optimal = {res.x[0]:.2f} MW, P2 Optimal = {res.x[1]:.2f} MW")
print(f"Total Biaya Minimum = ${res.fun:.2f} / jam")
\end{lstlisting}
\end{frame}

% --- SECTION 4 ---
\section{Kesimpulan}

\begin{frame}{Rangkuman dan Diskusi Pertemuan 1}
    \begin{tcolorbox}[colback=unibblue!10,colframe=unibblue,title=\textbf{Poin Utama Minggu ke-1}]
        \begin{enumerate}
            \item Pemodelan matematis presisi adalah langkah terpenting sebelum memilih metode optimisasi.
            \item Pemahaman klasifikasi masalah (LP, NLP, Convex, MIP) menentukan jenis *solver* dan algoritma yang tepat.
            \item Bahasa komputasi modern seperti Python dan Julia memudahkan implementasi langsung dari formulasi teori ke kode program.
        \end{enumerate}
    \end{tcolorbox}

    \vspace{0.4cm}
    \begin{center}
        \Large \textbf{Ada Pertanyaan / Diskusi?} \\
        \vspace{0.2cm}
        \small \textit{Lembar Kerja Mahasiswa (LKM 1) dan Problem Set 1 dapat diunduh pada portal perkuliahan.}
    \end{center}
\end{frame}

\end{document}
